\documentclass[11pt,a4paper]{article}
\usepackage[utf8]{inputenc}
\usepackage[T1]{fontenc}
\usepackage{amsmath,amssymb,amsthm}
\usepackage{booktabs}
\usepackage{graphicx}
\usepackage{microtype}
\usepackage{xcolor}
\usepackage{float}
\usepackage[margin=1.1in]{geometry}
\usepackage[colorlinks=true,linkcolor=blue!60!black,citecolor=blue!60!black,urlcolor=blue!60!black]{hyperref}
\graphicspath{{figures/}}

\newtheorem{theorem}{Theorem}[section]
\newtheorem{proposition}[theorem]{Proposition}
\newtheorem{corollary}[theorem]{Corollary}
\theoremstyle{definition}
\newtheorem{definition}[theorem]{Definition}
\theoremstyle{remark}
\newtheorem{remark}[theorem]{Remark}

\newcommand{\Z}{\mathbb{Z}}
\newcommand{\frc}{\operatorname{frac}}
\DeclareMathOperator{\dist}{dist}

\title{Reflecting Walls Are Self-Collisions:\\
  Cellular Automata in a Mirror Box}
\author{Miguel Esteban P\'erez Ibarra\\[2pt]
  \small Independent researcher, Mexico\\
  \small \href{mailto:supersluky@gmail.com}{\texttt{supersluky@gmail.com}}}
\date{\small Code and figures: \url{https://github.com/MiguelPerezI/mirror-walls}}

\begin{document}
\maketitle

\begin{abstract}
A cellular automaton with a reflection-symmetric local rule, run on a
box whose walls are implemented by mirror ghost cells, is exactly the
same automaton on the doubled periodic domain restricted to
mirror-symmetric configurations. This is the method of images, and it
has a dynamical consequence that we make the organizing principle of the
paper: a localized structure that reaches a wall collides head-on with
its own mirror image, at the one relative phase the wall admits, and at
a stratum of codimension $c$ of the box it meets $2^c-1$ images. What a
wall does to a structure is therefore read off from the bulk collision
table, restricted to symmetric collisions. We test this on three
automata. In Wolfram's continuous automaton (three-cell mean plus $0.398$
modulo $1$) the symmetric collision is elastic and the walls are
invisible to the population: bulk densities in the box agree with the
torus to three digits. In elementary rule~54 we show by exhaustive
enumeration that each direction of motion has a single phase class, so
every head-on collision of two single movers is mirror-symmetric and
annihilates; the wall is a perfect absorber, and because a symmetric
state carries as many left as right movers, a sparse box drains all of
its movers (in $90$ of $90$ launches) while a torus keeps the unbalanced
excess (in $46$ of $90$). In rule~110, which is not reflection
symmetric, the principle fails and a wall becomes a phase source that
can replace the background. In two and three dimensions the boundary
strata, indexed by the ternary cube, are quiet in proportion to their
codimension, and phase walls are cleared from a margin tens of cells deep
by annihilation with their images. We discuss in what limited sense the
box acts as a selecting environment.
\end{abstract}

\section{Introduction}\label{sec:intro}

Localized structures of cellular automata (gliders, particles, defects)
are usually studied on a ring or a torus, where nothing distinguishes
one place from another. Finite domains with reflecting walls are the
natural setting of the partial differential equations that many
automata discretize, and they are where the question of an
\emph{environment} first arises: a wall is a fixed feature that every
structure eventually meets. This paper asks what reflecting walls do to
localized structures, and answers it with one exact observation and
three case studies.

The observation (Section~\ref{sec:reflection}) is the method of images:
for a rule that is symmetric under $i\mapsto-i$, the box with mirror
ghosts is the periodic system on the doubled domain restricted to
mirror-symmetric states. It follows that a wall is not a new kind of
object for the dynamics. A structure meeting it undergoes the head-on
collision with its own mirror image, at the single relative phase that
puts the symmetry axis on the wall node. The effect of the wall is
therefore the outcome of one entry of the bulk collision table, and
whether the wall reflects, absorbs or transforms structures is decided
by the rule, not by the boundary. Edges and corners of a box are fixed
sets of several commuting reflections; we index them by the ternary
cube \cite{TernaryCube}, in which each node of an $n$-dimensional box
carries a word in $\{0,1,2\}^n$ recording, per axis, whether it lies on
the lower wall, in the interior, or on the upper wall.

The case studies are chosen so that the symmetric collision has three
different outcomes. In the one-dimensional continuous automaton of
Wolfram \cite[ch.~4]{NKS} at $a=0.398$ it is elastic, and the walls are
invisible (Section~\ref{sec:cca}). In elementary rule~54
\cite{Boccara,HansonCrutchfield,Martinez2006,Martinez2014} it annihilates;
we show that in this rule it is the \emph{only} head-on collision of two
single movers, so walls and bulk destroy movers by the same event, and
the difference between box and torus is a conservation argument: a
symmetric state is balanced (Section~\ref{sec:r54}). Rule~110
\cite{Cook} is not reflection symmetric, and its walls behave as sources
(Section~\ref{sec:r110}). In two and three dimensions
(Section~\ref{sec:2d}) we measure activity per boundary stratum for the
continuous automaton and find quiet strata, quieter with increasing
codimension, and a deep boundary layer where the bulk is organized in
phase domains.

\paragraph{What is and is not new.} The method of images is classical,
and the glider collisions of rules~54 and~110 are catalogued
\cite{Boccara,HansonCrutchfield,Martinez2006,Martinez2014}. What we add
is the reading of reflecting walls as symmetric self-collisions, the
balance argument that turns it into a statement about populations, an
exhaustive enumeration of the phase classes of the rule-54 movers in
the uniform background, and measurements of how boundary strata of every
codimension affect a two- and three-dimensional automaton. We found no
prior study of localized structures of these automata against
reflecting walls.

\section{Setting}\label{sec:setting}

\paragraph{Continuous automata.} On an $n$-dimensional grid with nodes
$\mathbf i$ and values $u_{\mathbf i}\in[0,1)$, the rule is
\begin{equation}\label{eq:cca}
  u^{l+1}_{\mathbf i}=\frc\Bigl(u^l_{\mathbf i}+\sum_{m=1}^nr_m\,\tilde\delta^2_{m,\tau_m(\mathbf i)}u^l_{\mathbf i}+a\Bigr),
\end{equation}
with $r_m=1/(2n+1)$, that is, the mean of the $(2n+1)$-point
neighbourhood plus a constant $a$, modulo~$1$. The operator
$\tilde\delta^2_{m,\tau_m}$ is the per-axis second difference of the
ternary-cube rule \cite{TernaryCube}: on an interior digit $\tau_m=1$ it
is the central difference, on a wall digit it uses the mirror ghost
value $u_{-1}:=u_{1}$ (lower wall) or $u_{N}:=u_{N-2}$ (upper wall).
This is the explicit heat step with Neumann walls followed by a shift and
a fold, a coupled map lattice in the sense of \cite{Kaneko}. Wolfram
reports class-4 behaviour in one dimension for $a=0.398$ and $a=0.4$
\cite[p.~244]{NKS}. Values are rounded to nine decimals after each step,
which makes the dynamics a map on a finite state set and runs
reproducible across machines.

\paragraph{Elementary automata.} For an elementary rule $f$ on
$\{0,1\}^N$, the box uses the same mirror ghosts, $u_{-1}:=u_1$ and
$u_N:=u_{N-2}$. A rule is reflection symmetric if
$f(\ell,c,r)=f(r,c,\ell)$; rule~54 is, rule~110 is not (its mirror image
is rule~124).

\paragraph{Background and defects.} Rules~54 and~110 have a periodic
background, the ether. For rule~54 it is the tiling by $1000$, of
temporal period~$4$; for rule~110 the tiling by $11000100110111$, of
temporal period~$7$. Since the ether is known exactly, we display and
count structures by comparing the state with the pure ether evolved in
the same geometry: a cell is \emph{off-ether} if it differs from it. For
rule~54 the mirror box is ether-compatible when its length is $4k+1$,
which puts both wall nodes on a~$1$ of the tile.

\section{The reflection principle}\label{sec:reflection}

Write $\mathcal R_m$ for the even reflection of a configuration across
the wall nodes of axis $m$, which turns the box grid
$\{0,\dots,N_m-1\}$ into the periodic grid of length $2(N_m-1)$, and
$\mathcal R=\mathcal R_1\cdots\mathcal R_n$ for the reflection into the
torus $\mathbb T=\prod_m\Z/2(N_m-1)\Z$.

\begin{theorem}[reflection principle]\label{thm:reflection}
Let $F$ be a local rule invariant under each axis reflection
$i_m\mapsto-i_m$ (for instance a symmetric per-axis stencil followed by
any pointwise map, or a reflection-symmetric elementary rule). Let $F_{\rm box}$ be $F$ on
the box with mirror ghosts and $F_{\rm per}$ the same rule on
$\mathbb T$. Then
\[
  \mathcal R\circ F_{\rm box}=F_{\rm per}\circ\mathcal R ,
\]
and $\mathcal R$ maps the box configurations bijectively onto the
configurations of $\mathbb T$ that are symmetric under all $n$
reflections, which form an invariant set of $F_{\rm per}$. This covers
\eqref{eq:cca}, its weighted variants \cite[p.~244]{NKS}, and every
reflection-symmetric elementary rule.
\end{theorem}

\begin{proof}
At a node with an interior digit the two rules read the same neighbours.
At a wall node the mirror ghost is exactly the value that the reflected
configuration carries at the position beyond the wall, so the stencils
agree. Symmetry of the stencil makes the update of a symmetric
configuration at $-i$ equal to its update at $i$, so symmetry is
preserved. Pointwise maps commute with both.
\end{proof}

This is the method of images, and for linear stencils it is the reason
Neumann problems can be solved by even extension. Numerically the
identity holds to the last digit: for \eqref{eq:cca} on $400$ cells for
$3000$ steps in one dimension and on $64^2$ cells for $1000$ steps in
two dimensions, the maximum difference between the box run and the
doubled periodic run is $0$ (script \texttt{reflection\_principle.py}).

\begin{corollary}[walls are self-collisions]\label{cor:self}
Under the hypotheses of Theorem~\ref{thm:reflection}, the evolution of
a structure that reaches a wall of the box is the evolution, in the
periodic system, of the head-on collision of that structure with its
mirror image, at the relative phase that puts the symmetry axis on the
wall node. At a node of codimension $c$ (a word of the ternary cube with
$c$ wall digits) the structure meets $2^c-1$ images under $c$ commuting
reflections. A configuration supported on a stratum is a configuration
of the periodic system that is symmetric under the corresponding
reflections.
\end{corollary}

Two consequences are used below. First, the wall realizes only the
\emph{symmetric} entries of the collision table, and it realizes them
for every structure that reaches it: if a rule has several head-on
phases with different outcomes, the wall picks out one, always the same.
Second, a stratum is quiet for a kinematic reason: the wall node is a
symmetry centre, so only mirror-symmetric local patterns can sit on it,
and a node of codimension $c$ has $c$ of its neighbour pairs identified,
so it has fewer independent neighbours to differ from.

\begin{proposition}[balance]\label{prop:balance}
Let a reflection-symmetric one-dimensional rule have, in a background,
movers of speed $+v$ and $-v$ that are mirror images of each other.
\begin{enumerate}
\item In the doubled system of a box, the number of right movers equals the number of left movers at all times.
\item Suppose that in some regime movers are neither created nor converted, that two movers of the same direction never meet, and that every head-on meeting of a right and a left mover annihilates both. Then on a ring the number $N_+-N_-$ is conserved and $|N_+-N_-|$ movers survive forever, while in the box every mover is eventually annihilated.
\end{enumerate}
\end{proposition}

\begin{proof}
(1) The reflection maps right movers to left movers and the doubled
state is symmetric. (2) Each annihilation removes one mover of each
direction, which conserves $N_+-N_-$. On a ring of length $L$ a right and
a left mover approach at relative speed $2v$ and meet within $L/(2v)$
steps, so movers are annihilated in pairs until only one direction
remains, and movers of one direction never meet. In the box, the
doubled ring has $N_+=N_-$ by (1), so all movers are annihilated.
\end{proof}

\section{Elastic walls: the continuous automaton}\label{sec:cca}

We take \eqref{eq:cca} in one dimension with $a=0.398$ from random data
and the activity indicator
$\chi_i=[\dist(u_i-u_{i+1},\Z)>0.05]$, which is Wolfram's difference
display thresholded. Figure~\ref{fig:class4_1d} shows the space-time
diagrams on a ring and in a box.

\begin{figure}[tbp]
\centering
\includegraphics[height=0.45\textheight]{class4_1d_a0.398_periodic.pdf}\hfill
\includegraphics[height=0.45\textheight]{class4_1d_a0.398_mirror.pdf}\hfill
\includegraphics[height=0.45\textheight]{class4_1d_wallzoom.pdf}
\caption{Class~4 in one dimension: $a=0.398$, three-cell mean plus $a$
modulo $1$, $300$ cells, $600$ steps downward from random data, same
seed, difference display (black where neighbours differ). Left: ring.
Middle: mirror walls at both sides. Right: the $60$ cells nearest the
left wall during steps $300$--$600$, magnified. Localized structures
move through a periodic background and at the wall return as their
mirror images, forming V shapes.}
\label{fig:class4_1d}
\end{figure}

\paragraph{The zoo.} Tracking clusters of active cells on a $600$-cell
ring over $3000$ steps gives one dominant mover, of speed $0.1$ cells per
step in either direction and width about $3.5$ cells, and a rare slow,
wider class (about $5\%$ of tracks), which is born in collisions and
lives about $60$ steps against about $105$ for the movers. A head-on
collision of two movers returns the pair, so by
Corollary~\ref{cor:self} a mover hitting a wall should come back, and
it does (Figure~\ref{fig:class4_1d}, right).

\paragraph{The walls are invisible to the population.} Over ten seeds,
$400$ cells and $10^4$ steps after a transient, counting active clusters
in the bulk (cells $50$ to $350$) at every step, the ring and the box
give $0.0723$ clusters per cell per step both, and fractions $0.406$ and
$0.408$ of wide clusters (width at least $6$). Mean lifetimes of both
classes agree between the geometries ($104$ against $105$ steps for
movers, $59$ against $60$ for the slow class). Total cluster counts over
$2\cdot10^4$ steps fluctuate around $30$ in both geometries with no
trend. Near the walls, the wall node has activity $0.09$ and its
neighbour $0.23$ against a bulk value of $0.28$, after which the bulk
value is reached: a quiet layer two cells deep, as
Corollary~\ref{cor:self} predicts, and nothing else. Walls that reflect
elastically do not select.

\section{Absorbing walls: rule 54}\label{sec:r54}

Rule~54 has the ether $1000$, movers $w_\pm$ of speed $\pm1$, stationary
breathers, and glider guns \cite{Boccara,HansonCrutchfield,Martinez2006,Martinez2014}.
Its head-on collisions are known to depend on the relative phase:
gliders ``in the same phase'' annihilate, gliders in anti-phase fuse into
a stationary breather \cite{Martinez2006}.

\paragraph{One phase class per direction.} We insert every binary word
of length up to $10$ at each of the four offsets into the ether, evolve
$160$ steps on a ring, and keep the outcomes that consist of a single
mover, identified as off-ether cells moving rigidly at speed one
(script \texttt{rule54\_exact.py}). Up to translation by the ether
period, exactly one shape is found for each direction: a pair of
adjacent off-ether cells, starting at residue $2$ modulo $4$ (right
movers) or $1$ modulo $4$ (left movers) at times divisible by~$4$; the
other outcomes are trains of two or more movers four cells apart. A
survey of every mover in $1500$ random launches of $20$ cells
(script \texttt{rule54\_phases.py}), $3\times10^5$ observations including
movers born in multi-body events, finds no other residue. So in the
uniform ether the phase of a mover is not a free parameter, and the
relative phase of a head-on collision of two single movers is fixed.

\paragraph{The collision and the wall event.} Placing a single right
and a single left mover at every admissible pair of offsets on a ring,
we find one collision class: it is mirror-symmetric about a node at all
times and annihilates both movers, leaving pure ether. This is the
``same phase'' annihilation of \cite{Martinez2006}; the anti-phase
fusion does not arise between single movers in the uniform ether. By
Corollary~\ref{cor:self} the wall event is the same collision, and
directly: a single right mover started at any offset in a box of length
$4k+1$ meets the wall in one way, and nothing is left. A train of four
movers meeting a single mover loses one member and continues; trains of
two and three meeting the wall lose all their members
(Figure~\ref{fig:r54events}).

\begin{figure}[tbp]
\centering
\includegraphics[width=0.9\textwidth]{rule54_events.pdf}
\caption{Rule~54, off-ether display (black where the state differs from
the exact ether), time downward, $160$ steps by $160$ cells. Left: a
single right and a single left mover meet on the ring and annihilate.
Middle: a single right mover reaches the mirror wall (right edge of the
panel) and vanishes: the same event, with the left mover supplied by
the reflection. Right: a single right mover meets a train of four left
movers; one pair annihilates and three continue.}
\label{fig:r54events}
\end{figure}

\paragraph{Sparse populations.} One random $8$-cell patch is placed at
the centre of the ether, in a ring of $400$ cells and in a box of $401$
cells with the same patch, and the movers are counted exactly (a right
mover is a pair of off-ether cells at residue $2$ that reappears four
cells to the right four steps later; script \texttt{rule54\_sparse.py}).
Of $200$ patches, $110$ produce a glider gun, which fills either
geometry with movers at the same rate; Table~\ref{tab:r54} summarizes
the other $90$.

\begin{table}[tbp]
\centering
\begin{tabular}{@{}lrrrrrr@{}}
\toprule
& \multicolumn{6}{c}{mean number of movers at step} \\
geometry & $100$ & $300$ & $600$ & $1000$ & $1500$ & $2000$ \\
\midrule
ring ($400$ cells) & 2.06 & 1.17 & 1.14 & 1.14 & 1.14 & 1.14 \\
box ($401$ cells)  & 2.06 & 0.00 & 0.00 & 0.00 & 0.00 & 0.00 \\
\bottomrule
\end{tabular}
\caption{Rule~54, $90$ sparse launches without guns, same patches in
both geometries. At step $2000$ movers survive in $46$ of the $90$ rings
and in none of the boxes. In every ring the survivors move in a single
direction, and $N_+-N_-$ at step $300$ equals $N_+-N_-$ at step $2000$ in
$88$ of the $90$ rings.}
\label{tab:r54}
\end{table}

This is Proposition~\ref{prop:balance} at work. The hypotheses hold in
the sparse regime by the enumeration above: every head-on meeting of
single movers annihilates, and a train loses one member for each single mover it
meets. On
the ring the unbalanced movers never meet a partner and live forever;
in the box each mover's partner is its own image, so all of them die,
each within one crossing of the box (Figure~\ref{fig:r54sparse}). In the
two exceptional rings the patch leaves a small stationary structure
behind, and the single mover is absorbed by it on its first pass,
an interaction outside the hypotheses of the proposition (stationary
structures can absorb, reflect or convert movers \cite{Martinez2006}).

\begin{figure}[tbp]
\centering
\includegraphics[width=0.85\textwidth]{rule54_sparse.pdf}
\caption{One sparse launch, off-ether display, $1200$ steps (every
second step shown), ring of $400$ cells (left) and box of $401$ cells
(right), same patch. The patch emits a train of four left movers and
one right mover. On the ring the right mover meets the train across the
periodic seam and annihilates one member; the other three circulate
forever. In the box every mover meets its image at a wall.}
\label{fig:r54sparse}
\end{figure}

\paragraph{Dense populations.} From random data on $1200$ cells the
population is sustained by guns, and walls and bulk both destroy movers
by the same collision. Over three seeds and six sampling times between
steps $1000$ and $6000$ the ring carries $107\pm9$ movers and the box
$94\pm13$: no difference beyond the fluctuations, which are dominated
by the birth and death of guns. The walls are perfect absorbers, but in a
dense population the bulk is nearly as effective.

\section{Asymmetric rules: rule 110}\label{sec:r110}

Rule~110 is not reflection symmetric, so Theorem~\ref{thm:reflection}
does not apply and a mirror wall is a genuinely new boundary rule, not a
self-collision. Its ether has spatial period $14$ and is not
mirror-compatible: the right wall, which is nearly ether-compatible,
hosts stationary structures and little else, but the left wall can be
kicked by an arriving glider into a state that grows a different periodic
background. That background invades the box behind a front, emitting
fast right-moving gliders as it advances (Figure~\ref{fig:r110}). Once
nucleated the invasion is total: in one of two $1400$-cell runs the
ether fraction fell from $0.90$ at step $5000$ to $0$ at step $8000$; in
the other it stayed at $0.96$ through step $8000$. The wall changes the
population by emission and eventually destroys the medium the gliders
live in. Without the reflection symmetry there is no principle that
predicts what a wall does, and the outcome depends on the details of
the boundary rule.

\begin{figure}[tbp]
\centering
\includegraphics[width=\textwidth]{rule110_defects_fixed.pdf}
\caption{Rule~110, defect display (black where the $14$-cell word is not
an ether word), time downward. Left: ring. Right: box with mirror
walls, same initial data; the left wall has nucleated a different
background, which advances into the box and emits gliders.}
\label{fig:r110}
\end{figure}

\section{Two and three dimensions}\label{sec:2d}

\paragraph{Regimes.} In two dimensions \eqref{eq:cca} is the five-point
mean plus $a$ modulo~$1$. A scan of $a$ in steps of $0.01$ on a $128^2$
torus from random data (script \texttt{scan2d.py}) gives a late activity
fraction that varies continuously with $a$, from $0.29$ near $a=0.01$ to
$0.90$--$0.92$ on the plateaus $a\in[0.34,0.43]$ and $[0.57,0.66]$, with dips
to $0.42$--$0.45$ near $a=\frac13,\frac12,\frac23$. There are no sharp
class-4 windows of the one-dimensional kind, and the one-dimensional
class-4 value $a=0.398$ is fully chaotic in two dimensions. Space-time
cuts (Figure~\ref{fig:cuts}) show three regimes: near $a=\frac13$ and
$\frac23$ the background is a period-three cycle and the dynamics
consists of slowly wandering \emph{phase walls} between its three
phases; near $a=\frac12$ there are intermittent short-lived spots on a
period-two background, in the manner of spatiotemporal intermittency
\cite{ChateManneville}; on the plateaus, spatiotemporal chaos. We found
no gliders in this family in two dimensions; its localized structures
are the phase walls and the spots.

\begin{figure}[tbp]
\centering
\includegraphics[width=\textwidth]{cuts2d_spacetime.pdf}
\caption{Space-time cuts (middle row of a $128^2$ torus against the last
$256$ of $1200$ steps, difference display) for
$a=0.33,\,0.49,\,0.51,\,0.67,\,0.02,\,0.398$. Wandering phase walls at
$0.33$ and $0.67$, intermittent spots at $0.49$ and $0.51$, a regular
lattice at $0.02$, chaos at $0.398$.}
\label{fig:cuts}
\end{figure}

\begin{figure}[tbp]
\centering
\includegraphics[width=\textwidth]{filmstrip2d_a0.33.pdf}\\[4pt]
\includegraphics[width=\textwidth]{filmstrip2d_a0.49.pdf}
\caption{The $128^2$ box with mirror walls, difference display, at
steps $0,50,100,200$ (top row of each strip) and $400,600,800,1200$
(bottom row). Upper strip, $a=0.33$: random data collapse within fifty
steps into a period-three background crossed by phase walls, which then
wander, join and pinch off; walls parallel to a side are cleared from a
margin several cells wide (tens of cells on larger boxes), walls perpendicular to it end on it. Lower
strip, $a=0.49$: intermittent spots on a period-two background, born and
dying everywhere except at the corners.}
\label{fig:box033}
\end{figure}

\paragraph{Activity per stratum.} We run the box with mirror walls,
classify every node by the codimension of its stratum ($0$ interior, $1$
face, $2$ edge or two-dimensional corner, $3$ three-dimensional corner),
and record the time-averaged activity indicator, defined as the maximum
over all existing neighbour differences of
$\dist(u_{\mathbf i}-u_{\mathbf j},\Z)>0.05$. Grids are $128^2$ for
$1500$ steps and $32^3$ for $800$ steps, averaged after a transient of
one third of the run, over eight seeds (script \texttt{strata\_seeds.py}).
Table~\ref{tab:strata} gives the results. A CuPy implementation of the
same automaton (\texttt{contca\_gpu.py}, \texttt{strata\_gpu.py})
reproduces every entry bit for bit and computes a row of the table in a
fraction of a second instead of about a minute on one CPU core, which
allows larger boxes: Table~\ref{tab:strata_big} repeats the
measurement on $1024^2$ for $3000$ steps and on $128^3$ for $1500$ steps,
with sixteen seeds.

\begin{table}[tbp]
\centering
\begin{tabular}{@{}llllll@{}}
\toprule
regime & grid & codim $0$ & codim $1$ & codim $2$ & codim $3$ \\
\midrule
chaos, $a=0.398$ & $128^2$ & 0.989(0) & 0.948(0) & 0.738(3) & \\
chaos, $a=0.398$ & $32^3$  & 0.997(0) & 0.993(0) & 0.981(0) & 0.922(3) \\
intermittent, $a=0.49$ & $128^2$ & 0.558(2) & 0.494(2) & 0.196(4) & \\
intermittent, $a=0.49$ & $32^3$  & 0.484(5) & 0.505(5) & 0.506(3) & 0.305(8) \\
phase walls, $a=0.33$ & $128^2$ & 0.510(5) & 0.340(3) & 0.172(11) & \\
chaos (3-D rule), $a=0.33$ & $32^3$ & 0.996(2) & 0.992(5) & 0.973(7) & 0.860(10) \\
\bottomrule
\end{tabular}
\caption{Time-averaged activity fraction per boundary stratum of the box
with mirror walls, mean over eight seeds with the standard error in the
last digits in parentheses.}
\label{tab:strata}
\end{table}

\begin{table}[tbp]
\centering
\begin{tabular}{@{}llllll@{}}
\toprule
regime & grid & codim $0$ & codim $1$ & codim $2$ & codim $3$ \\
\midrule
chaos, $a=0.398$ & $1024^2$ & 0.989(0) & 0.949(0) & 0.745(3) & \\
chaos, $a=0.398$ & $128^3$  & 0.997(0) & 0.993(0) & 0.981(0) & 0.922(1) \\
intermittent, $a=0.49$ & $1024^2$ & 0.551(0) & 0.494(0) & 0.207(2) & \\
intermittent, $a=0.49$ & $128^3$  & 0.451(0) & 0.474(0) & 0.487(1) & 0.309(2) \\
phase walls, $a=0.33$ & $1024^2$ & 0.561(1) & 0.351(1) & 0.165(6) & \\
chaos (3-D rule), $a=0.33$ & $128^3$ & 0.999(0) & 0.999(0) & 0.995(0) & 0.898(10) \\
\bottomrule
\end{tabular}
\caption{The measurement of Table~\ref{tab:strata} on larger boxes:
$1024^2$ for $3000$ steps and $128^3$ for $1500$ steps, sixteen seeds,
same transient, threshold and rounding (script \texttt{strata\_gpu.py}).
The strata with many nodes now have standard errors of at most $0.001$; the
corner strata keep four or eight nodes per run and do not sharpen with
the box size.}
\label{tab:strata_big}
\end{table}

Four observations. First, in the chaotic regime the activity decreases
monotonically with the codimension of the stratum, in both dimensions,
and the boundary layer is one cell deep: at distance one from the
nearest wall the two-dimensional activity is already $0.982$ against
$0.989$ in the bulk. In three dimensions the values are the same to
three decimals from $32^3$ to $128^3$; in two dimensions they change by
at most $0.007$ from $128^2$ to $1024^2$. This is the kinematic effect of
Corollary~\ref{cor:self}: a node of codimension $c$ has $c$ of its
neighbour pairs identified by symmetry, so it has fewer independent
neighbours to differ from, and the effect compounds with $c$. Part of
the trend is therefore a property of the indicator, which takes a
maximum over fewer independent differences on the strata. Second, in the
intermittent regime the corner is strongly suppressed in both
dimensions ($0.20$ against $0.56$ in two dimensions, $0.31$ against
$0.48$ in three on $32^3$ and against $0.45$ on $128^3$), while the
three-dimensional faces and edges sit above the bulk value: the faces
by $0.02$ on both boxes, the edges by $0.02$ on $32^3$ and $0.04$ on
$128^3$.
In this regime the wall is felt beyond one cell, in the opposite sense:
the activity overshoots the bulk at distances $2$--$4$ in two dimensions
($0.59$ against $0.55$) and at distance $1$ in three ($0.49$ against
$0.45$), and relaxes within five to seven cells. So in two dimensions the spots avoid the
corners but not the walls, and in three dimensions only the eight corner
nodes are special; the twelve edges, which are the analogue of the
two-dimensional corners, are not. Third, the phase-wall regime is the
one in which the box is felt at a distance: the activity profile
against the distance to the nearest wall is
$0.34,\,0.38,\,0.43,\,0.45,\,0.46,\,0.47$ for distances $0$ to $5$ on
$128^2$, still rising at distance five. On $1024^2$ the profile keeps
rising, to $0.50$ at distance $7$, $0.53$ at $16$ and $0.55$ at $24$, and
reaches the bulk value $0.56$ only $35$--$40$ cells from the wall. So the
codimension-$0$ entry of Table~\ref{tab:strata} for this regime is not a
bulk value but an average over a box most of which lies in the margin
($0.51$ on $128^2$, $0.54$ on $256^2$, $0.55$ on $512^2$, $0.56$ on
$1024^2$), while the face and corner entries are stable in the box size.
Corollary~\ref{cor:self} says why: a phase wall parallel to a mirror
meets its image, a phase wall of opposite orientation, and the pair
annihilates when the strip between them collapses, so a margin along
each mirror is cleared of parallel phase walls; phase walls
perpendicular to the mirror are their own images, end on it, and are
the activity that remains at distance zero. This is the two-dimensional
counterpart of the absorbing walls of rule~54. Fourth, the same
$a=0.33$ is chaotic under the seven-point rule in three dimensions, with
the one-cell layer of the chaotic regime (on $64^3$, $96^3$ and $128^3$
the faces equal the bulk, $0.999$, and only the edges, $0.994$--$0.995$,
and the corners, $0.89$--$0.90$, depart), which shows that the regime, not the
constant, is what the strata respond to.

\section{Discussion: is the box a selecting environment?}\label{sec:disc}

Wong et al.\ \cite{Wong2023} have proposed that the functional
information of a system, in the sense of
\cite{Szostak2003,Hazen2007}, increases when many configurations are
subjected to selection for some function, and have asked which
environments act as selectors. A box with mirror walls is about the
simplest fixed environment a cellular automaton can have, and the
results above say precisely what it selects. The ``function'' is
surviving one's own reflection: the wall applies to every structure that
reaches it the symmetric entry of the collision table, and keeps exactly
the structures for which that entry is elastic. In the continuous
automaton of Section~\ref{sec:cca} every structure passes and the box
selects nothing. In rule~54 no mover passes, so the box filters a sparse
population to its stationary members, while the ring keeps a conserved
excess of movers; this is differential persistence imposed by the
environment, and it is the only regime in which we see the box change a
population. It is one-shot filtering toward a fixed survivor set, not
the open-ended accumulation the law of \cite{Wong2023} describes:
nothing is copied, and nothing varies. In rule~110 the walls do not
filter but emit.

Selection among structures, as opposed to against all of one kind,
requires a reflection-symmetric rule with several mover species whose
symmetric self-collisions differ, some elastic and some not. Rule~54 has
the symmetry but one mover species; rule~110 has a rich zoo
\cite{Cook} but not the symmetry. Two-dimensional continuous automata of
Lenia type \cite{Lenia} have both, and add the corner, a four-body
symmetric collision, as a second and independent filter. The experiment
of this paper (exhaustive wall events, sparse launches with the same
seeds in both geometries, exact counting) transfers to them unchanged.

\section*{Code}\sloppy
All scripts, the figures and the source of this article are public at
\url{https://github.com/MiguelPerezI/mirror-walls}. In \texttt{code/}:
\texttt{contca.py} (rule \eqref{eq:cca} with ternary-cube walls) and
\texttt{reflection\_principle.py} (Theorem~\ref{thm:reflection});
\texttt{view1d.py}, \texttt{tracker.py}, \texttt{catalogue2.py},
\texttt{slowclass2.py}, \texttt{perstep.py}, \texttt{population.py} and
\texttt{wall1d.py} (Section~\ref{sec:cca}); \texttt{eca.py},
\texttt{rule54\_exact.py}, \texttt{rule54\_phases.py},
\texttt{rule54\_sparse.py}, \texttt{rule54\_charge.py} and
\texttt{fig\_rule54.py} (Section~\ref{sec:r54}); \texttt{eca\_box.py},
\texttt{rule110\_ether.py} and \texttt{fig\_rule110.py}
(Section~\ref{sec:r110}); \texttt{scan2d.py}, \texttt{cuts2d.py},
\texttt{filmstrip2d.py} and \texttt{strata\_seeds.py}
(Section~\ref{sec:2d}); \texttt{contca\_gpu.py} and
\texttt{strata\_gpu.py} are CuPy versions of the automaton and of the
strata measurement, bit-identical to the NumPy ones
(Table~\ref{tab:strata_big}). The script \texttt{reproduce.sh} at the top
of the repository runs all of them with the arguments used here, writes
their outputs to \texttt{results/} and regenerates the figures, and
\texttt{check\_results.py} compares every number quoted in this article
with those outputs. \texttt{view2d.py} is a live viewer of the
two-dimensional box (\texttt{python3 view2d.py 0.33 160 mirror}).

\begin{thebibliography}{99}
\bibitem{TernaryCube}
M.~E.~P\'erez Ibarra,
\emph{One case, not twenty-seven: the ternary cube as the index set for
boundary stencils of finite-difference schemes in any dimension},
in preparation (2026).
\bibitem{NKS}
S.~Wolfram, \emph{A New Kind of Science}, Wolfram Media, 2002.
\bibitem{Kaneko}
K.~Kaneko, \emph{Overview of coupled map lattices}, Chaos 2 (1992) 279--282.
\bibitem{Boccara}
N.~Boccara, J.~Nasser, M.~Roger, \emph{Particlelike structures and their
interactions in spatiotemporal patterns generated by one-dimensional
deterministic cellular-automaton rules}, Phys.\ Rev.\ A 44 (1991)
866--875.
\bibitem{HansonCrutchfield}
J.~E.~Hanson, J.~P.~Crutchfield, \emph{Computational mechanics of
cellular automata: an example}, Physica D 103 (1997) 169--189.
\bibitem{Martinez2006}
G.~J.~Mart\'inez, A.~Adamatzky, H.~V.~McIntosh, \emph{Phenomenology of
glider collisions in cellular automaton Rule~54 and associated logical
gates}, Chaos Solitons Fractals 28 (2006) 100--111.
\bibitem{Martinez2014}
G.~J.~Mart\'inez, A.~Adamatzky, H.~V.~McIntosh, \emph{Complete
characterization of structure of Rule~54}, Complex Systems 23 (2014);
arXiv:1410.3096.
\bibitem{Cook}
M.~Cook, \emph{Universality in elementary cellular automata}, Complex
Systems 15 (2004) 1--40.
\bibitem{ChateManneville}
H.~Chat\'e, P.~Manneville, \emph{Transition to turbulence via
spatio-temporal intermittency}, Phys.\ Rev.\ Lett.\ 58 (1987) 112--115.
\bibitem{Lenia}
B.~W.-C.~Chan, \emph{Lenia and expanded universe}, arXiv:2005.03742.
\bibitem{Wong2023}
M.~L.~Wong, C.~E.~Cleland, D.~Arend~Jr., S.~Bartlett, H.~J.~Cleaves~II,
H.~Demarest, A.~Prabhu, J.~I.~Lunine, R.~M.~Hazen, \emph{On the roles of
function and selection in evolving systems}, Proc.\ Natl.\ Acad.\ Sci.\
USA 120 (2023) e2310223120.
\bibitem{Szostak2003}
J.~W.~Szostak, \emph{Functional information: molecular messages}, Nature
423 (2003) 689.
\bibitem{Hazen2007}
R.~M.~Hazen, P.~L.~Griffin, J.~M.~Carothers, J.~W.~Szostak,
\emph{Functional information and the emergence of biocomplexity},
Proc.\ Natl.\ Acad.\ Sci.\ USA 104 (2007) 8574--8581.
\end{thebibliography}

\end{document}
